% Chapter 7, generated by make_chapters.py from papers/Ch07_Bio_08_DS_Lattice_Sum/anccft21.tex.
% Source paper: Bio 8
\chapter[Double-stranded assembly as a lattice sum]{Double-stranded genome assembly as a lattice sum, and the inversion-flip group}\label{chap:ch07}
\begin{chaptersummary}
The two companion notes on double strands ended with a question: the sandpile group
$\K(D_k)$ of the reverse-complement double cover counts Eulerian circuits through both
strands at once, its deck involution is an anti-automorphism carrying a symmetric linking
form and no metabolizer, and what a double-stranded reconstruction count is built from was
left open. This note answers it. A circular sequence $T$ is a double-stranded reconstruction
of $S$ at word length $k$ when $\spec(T)+\spec(\bar T)=\spec(S)+\spec(\bar S)$; writing $m$
for the arc multiplicities of $D_k$ and $c=\spec(T)$, this says $c+\rho c=m$, and with
$f=2c-m$ it says that $f$ lies in the $(-1)$-eigenlattice $\Lm$ of $\rho$ on the circulation
lattice of $D_k$, in the box $|f|\le m$, in the parity class of $m$. The count is therefore a
\emph{lattice sum}, $\#\DS(S,k)=\sum_f \Nseq\big((m+f)/2\big)$, each term the BEST count of a
multigraph divided by $\prod_a c(a)!$ --- and that division, invisible on one strand where it
is a constant, is mandatory here because it varies from point to point. The rank of $\Lm$ is
$(|A|-|V|+2-\#\mathrm{fix}_A-\#\mathrm{fix}_V)/2$ by the Hopf trace formula applied to an
arc-reversing involution; it is \emph{not} the number of inverted repeats, and the guess that
it is --- made in the planning of this note --- is refuted at $k=12$, where one inverted-repeat
pair and one palindrome meet a lattice of rank four. For $\phi$X174 the sum gives $10$
reconstructions at $k=12$ and $864$ at $k=11$, against $|\K(D_k)|=33$ and $25\,318\,912$, and
an independent enumeration of closed walks with the strand-pairing constraint returns the
same numbers point by point. The second half concerns moves. A block $w\,X\,\bar w$ of $T$ may
be replaced by $w\,\overline{X}\,\bar w$; every such flip is a flip between a visit of a
vertex $v$ of $D_k$ and a later visit of $\rho(v)$, and it changes the spectrum by
$-c_B+\rho c_B$. For $\phi$X174 the $864$ reconstructions at $k=11$ form a \emph{single}
component under flips alone --- no flip fixes a spectrum, so the same-spectrum repeat
exchanges of the single-strand theory are compositions of flips through other classes, of
length four, six or eight --- and the five inverted-repeat pairs of $S$ do not suffice: six of
the $218$ spectrum classes need a flip at a pair $\{a,\bar a\}$ with $a$ a single-strand
repeat, available only after an earlier flip has imported $\bar a$. Connectivity is a property
of the genome and not a theorem: a sequence with repeats and no inverted repeat has a discrete
flip graph.
\end{chaptersummary}

\section{What was asked, and what is true}

The single-strand theory of these notes reads the reconstruction count of a circular sequence
$S$ at word length $k$ off the critical group of its de Bruijn multidigraph $G_k(S)$ through
the BEST theorem \cite{BEST,SmithTutte}. The first companion note Chapter~\ref{chap:ch05} passed to the
reverse-complement double cover $D_k(S)$, the multidigraph whose arcs are the $(k{+}1)$-mer
occurrences in $S$ and in $\bar S$, and found that everything single-stranded applies to it
verbatim but that its involution $\rho(v)=\bar v$ reverses arcs. The second Chapter~\ref{chap:ch06} proved
that this is fatal for every covering-theoretic route: $P\Delta P=\Delta^{\mathsf T}$, the
eigenlattices of $\rho$ are not $\Delta$-stable, the sandpile group $\K(D_k)$ sits in no
covering sequence, and what $\rho$ provides instead is a symmetric linking form on $\K(D_k)$
which for $\phi$X174 has no metabolizer because $|\K(D_k)|$ is not a square. That note closed
with the statement that whatever counts double-stranded reconstructions is not the order of a
distinguished subgroup of $\K(D_k)$, and left open what it is.

The answer is that it is not a group order at all. $\K(D_k)$ counts Eulerian circuits of
$D_k$, which are single closed walks through all $2N$ arcs, that is, through both strands.
A double-stranded reconstruction is a circular sequence $T$ of length $N$ whose
$(k{+}1)$-mer spectrum, together with that of $\bar T$, is the spectrum of $S$ together with
that of $\bar S$; it is a closed walk through \emph{half} the arcs, whose $\rho$-image walks
the other half. Writing $c$ for the spectrum of $T$ and $m$ for the multiplicities of $D_k$,
the condition is $c+\rho c=m$. That is a linear condition, and with $f=2c-m$ it becomes the
condition that $f$ be a $\rho$-anti-invariant circulation in the box $|f|\le m$ with the parity
of $m$. The reconstruction count is therefore a sum over the lattice points of a polytope of
the sequence count of the multigraph at each point (Theorem~\ref{ch07:thm:lattice}). The rank of
the lattice has a closed form (Theorem~\ref{ch07:thm:rank}), and the sum, evaluated for $\phi$X174,
gives $10$ and $864$ where $|\K(D_k)|$ gives $33$ and $25\,318\,912$ (Table~\ref{ch07:tab:lattice}).
An enumeration that never mentions the lattice --- closed walks of $D_k$ with the
strand-pairing constraint, bucketed by spectrum --- returns exactly the connected lattice
points, with exactly the predicted count at each (check (E)).

The second half of the note is about moves between reconstructions. If $T$ contains a $k$-mer
$w$ and, later, its reverse complement $\bar w$, the block between and including them may be
reverse-complemented in place; this is the inversion of classical genome rearrangement, and it
is the only move that changes which strand a stretch of $T$ is read from. For $\phi$X174 the
flip graph on all $864$ reconstructions at $k=11$ is connected (Theorem~\ref{ch07:thm:flips}),
and the way it is connected is informative: no flip ever preserves the spectrum, so the
same-spectrum moves of the single-strand theory --- exchanging the two copies of a repeat ---
arise here only as compositions of flips passing through other spectra, always of even length
at least four. This is the de Bruijn-graph shadow of the fact, familiar from sorting by
reversals \cite{HP99,BP96}, that a transposition is a product of reversals; here each reversal
must be flanked by an inverted-repeat pair, and the price is four rather than three.

Three corrections to the specification of this note are recorded in
Remark~\ref{ch07:rem:corrections}: the rank of $\Lm$ is not the number of inverted repeats; the
per-point count must be divided by $\prod c(a)!$; and flip-connectivity is a property of the
genome, not a theorem.

Throughout, $S$ is a circular sequence of length $N$ over $\{A,C,G,T\}$, $\bar S$ its reverse
complement, $\bar w$ that of a word $w$, and a \emph{circular sequence} is a word up to
rotation. All computations are exact; the scripts \texttt{ds\_lattice.py} and
\texttt{ds\_flips.py} reproduce every number below and print pass/fail batteries of nine and
four checks.

\section{The double cover, its contents, and the spectrum of a reconstruction}

\begin{definition}\label{ch07:def:cover}
$D_k=D_k(S)$ is the multidigraph of Chapter~\ref{chap:ch05}: one vertex per $k$-mer occurring at least
twice in $S$ and $\bar S$ together, one arc per occurrence of a $(k{+}1)$-mer in $S$ or in
$\bar S$, unitig-contracted so that an arc is a maximal segment between two consecutive
repeated $k$-mers on one strand. An arc is thus a \emph{labelled} object (a strand and a
position); its \emph{content} is the segment string from the start of its tail $k$-mer to the
end of its head $k$-mer; two arcs with the same content are parallel copies. Let $V$ be the
vertex set, $A$ the set of contents, $m\colon A\to\Z_{\ge1}$ the number of arcs of each
content, $\ell(a)$ the number of $(k{+}1)$-mers of the content $a$, and $\partial\colon
\Z^A\to\Z^V$, $\partial a=\mathrm{head}(a)-\mathrm{tail}(a)$. The involution $\rho$ sends a
vertex $w$ to $\bar w$, a content $a$ to $\bar a$, and the arc at position $p$ of one strand to
the arc at the mirrored position of the other.
\end{definition}

\begin{lemma}\label{ch07:lem:rho}
\begin{enumerate}
\item[\rm(a)] $\rho$ is a fixed-point-free involution on arcs, an involution on contents with
$m(\bar a)=m(a)$ and $\ell(\bar a)=\ell(a)$, and reverses direction:
$\partial(\bar a)=-\rho\,\partial a$.
\item[\rm(b)] $\rho$ is free on vertices if and only if no repeated $k$-mer is its own
reverse complement, which is automatic for odd $k$; it is free on contents if and only if no
content is its own reverse complement, and a content with $m(a)\ge2$ is a single
$(k{+}1)$-mer, so this is automatic for even $k$. On a fixed content $m(a)$ is even.
\item[\rm(c)] $D_k$ has one component if $S$ carries an inverted repeat of length at least
$k$ and two otherwise, the two being the $S$-half and the $\bar S$-half, exchanged by $\rho$.
\end{enumerate}
\end{lemma}

\begin{proof}
(a) An arc is an occurrence on a strand and $\rho$ exchanges the strands, so it is free on
arcs; its image on contents is reverse complementation, which preserves length and, being a
bijection on occurrences, preserves multiplicity. The arc $u\to v$ of content $a$ maps to the
arc $\bar v\to\bar u$ of content $\bar a$.
(b) The vertex statement is Chapter~\ref{chap:ch06}, Prop.~2. A content with two parallel copies is a
repeated segment, so its interior $k$-mers are repeated; but the interior $k$-mers of a
unitig segment are by construction not repeated, so the segment has no interior $k$-mer and
is a single $(k{+}1)$-mer, of odd length when $k$ is even, and no word of odd length equals
its reverse complement. If $a=\bar a$ then the occurrences of $a$ are paired by $\rho$ with
occurrences of $a$, in pairs, so $m(a)$ is even.
(c) is Chapter~\ref{chap:ch05}, Thm.~3: all $k$-mers of $S$ are joined by the closed walk $S$, all those of
$\bar S$ by $\bar S$, and the two walks meet exactly at $k$-mers common to $S$ and $\bar S$.
\end{proof}

Part (b) has a consequence worth stating: $\rho$ is never free on vertices and contents at
once, since palindromic $k$-mers need even $k$ and palindromic $(k{+}1)$-mers need odd $k$.
For $\phi$X174, $k=12$ has one palindromic repeated $12$-mer (\texttt{TGGAGGCCTCCA} at
$4483$) and no fixed content; $k=11$ has no fixed vertex and one fixed content, the $12$-mer
just named, now read as an arc between its two halves. The same word moves from one column
of Table~\ref{ch07:tab:rank} to the other as $k$ changes parity.

\begin{definition}\label{ch07:def:ds}
A \emph{double-stranded reconstruction} of $S$ at level $k$ is a circular sequence $T$ with
\[
  \spec(T)+\spec(\bar T)=\spec(S)+\spec(\bar S)
\]
as multisets of $(k{+}1)$-mers. $\#\DS(S,k)$ is their number.
\end{definition}

\begin{lemma}\label{ch07:lem:spectrum}
Let $T$ be a double-stranded reconstruction. Every $(k{+}1)$-mer of $T$ is an arc of $D_k$,
$T$ traverses each unitig segment of multiplicity one either once or not at all, and the
spectrum of $T$ is therefore a well-defined vector $c=\spec(T)\in\Z_{\ge0}^A$ with
$c+\rho c=m$ and $\partial c=0$. Moreover $|T|=N$.
\end{lemma}

\begin{proof}
The spectrum identity puts every $(k{+}1)$-mer of $T$ in the spectrum of $S$ or $\bar S$. A
non-repeated $k$-mer $w$ of $D_k$ has one arc in and one arc out; it occurs once in $S\bar S$,
hence once in $T\bar T$, hence at most once in $T$, and if it occurs in $T$ then the
$(k{+}1)$-mers of $T$ through it are its unique in- and out-arcs, so $T$ continues along the
unitig in both directions. This makes $c$ constant on unitigs. $c+\rho c=m$ is the spectrum
identity read on contents, since $\spec(\bar T)=\rho\,\spec(T)$; $\partial c=0$ because $T$ is
a closed walk. Finally $\sum_a c(a)\ell(a)+\sum_a c(\bar a)\ell(\bar a)=\sum_a m(a)\ell(a)=2N$
and the two sums on the left are equal.
\end{proof}

\section{The lattice sum}

\begin{definition}
$Z_1=\ker\partial\subset\Z^A$ is the circulation lattice of $D_k$, $\rho$ acts on it by
$(\rho f)(a)=f(\bar a)$, and $\Lp,\Lm$ are its $(\pm1)$-eigenlattices. For a circulation $c$
with $0\le c\le m$ let $G_c$ be the multigraph with $c(a)$ labelled parallel arcs of content
$a$, $t(c)$ its number of spanning arborescences converging on a root (independent of the
root, $G_c$ being balanced), $\Loc(c)=\prod_v(d^+_c(v)-1)!$ over the vertices of positive
$c$-degree, and
\[
  \Nseq(c)\;=\;\frac{t(c)\,\Loc(c)}{\prod_a c(a)!}\,.
\]
For a circular sequence $T$ let $d(T)$ be the number of rotations fixing $T$.
\end{definition}

\begin{theorem}[The lattice sum]\label{ch07:thm:lattice}
The map $T\mapsto\spec(T)$ is a surjection from the double-stranded reconstructions of $S$ at
level $k$ onto
\[
  P(S,k)=\{\,c\in\Z^A:\ 0\le c\le m,\ c+\rho c=m,\ \partial c=0,\ \supp c \text{ connected}\,\},
\]
and with $f=2c-m$ the set $P(S,k)$ is the set of $f\in\Lm$ with $|f|\le m$, $f\equiv m
\pmod 2$ and connected support of $(m+f)/2$. The fibre over $c$ has weighted cardinality
$\sum_{\spec T=c}1/d(T)=\Nseq(c)$. Hence
\[
  \sum_{T}\frac1{d(T)}\;=\;\sum_{f\in\Lm\cap\{|f|\le m,\ f\equiv m\ (2)\}}
  \Nseq\!\Big(\frac{m+f}{2}\Big),
\]
where the terms with disconnected support vanish, and when every reconstruction is aperiodic
the left side is $\#\DS(S,k)$.
\end{theorem}

\begin{proof}
Lemma~\ref{ch07:lem:spectrum} puts $\spec(T)$ in $P(S,k)$; the support is connected because $T$ is
one closed walk. Conversely a point $c\in P(S,k)$ defines a balanced connected multigraph
$G_c$, which has an Eulerian circuit, and reading the circuit gives a circular sequence $T$
with $\spec(T)=c$ and, by $c+\rho c=m$, with the spectrum identity. The substitution $f=2c-m$
is a bijection between $\{0\le c\le m\}$ and $\{|f|\le m,\ f\equiv m\ (2)\}$; $\partial c=0$
is $\partial f=0$ since $\partial m=0$; and $c+\rho c=m$ is $f+\rho f=0$.

For the fibre, the BEST theorem \cite{BEST,SmithTutte} counts the Eulerian circuits of
$G_c$, as cyclic sequences of its labelled arcs, by $t(c)\Loc(c)$. Forgetting the labels
sends a circuit to a circular sequence with spectrum $c$. Over a fixed $T$, a circuit is an
assignment of the $c(a)$ labelled copies of each content to the $c(a)$ occurrences of that
content around $T$, $\prod_a c(a)!$ assignments in all, and two assignments give the same
cyclic labelled sequence exactly when they differ by a rotation of $T$ onto itself. The
rotation group of $T$, of order $d(T)$, acts freely on assignments, so the fibre has
$\prod_a c(a)!/d(T)$ elements, and summing over $T$ gives
$t(c)\Loc(c)=\prod_a c(a)!\sum_T 1/d(T)$. If the support of $c$ is disconnected then $G_c$
has no spanning arborescence and $t(c)=0$.
\end{proof}

\begin{remark}[The normalisation is not optional]\label{ch07:rem:norm}
On one strand there is a single lattice point, $c=m$, and $\prod_a m(a)!$ is a constant; the
single-strand notes count in the \emph{arc-rooted} convention $t(m)\Loc(m)$ and never had to
divide. Across the lattice the divisor $\prod_a c(a)!$ changes from point to point, and the
un-normalised sum $\sum_f t(c_f)\Loc(c_f)$ is $\sum_T \prod_a c_T(a)!/d(T)$, which counts each
reconstruction with a weight of its own. The two agree whenever $m\le1$, which is the case at
$k=12$ for $\phi$X174, and disagree as soon as a content has two copies: at $k=11$ the
un-normalised sum is $2756$ and the count is $864$ (check (N)). The single-strand number
quoted for $\phi$X174 at $k=11$ in Chapter~\ref{chap:ch01}, $|\K(G_{11})|\Loc=132$, is the arc-rooted
count; the number of circular sequences with the spectrum of $S$ is $132/4=33$, and $33$ is
what appears at the point $c=\spec(S)$ of Table~\ref{ch07:tab:lattice}.
\end{remark}

\begin{remark}[Where the linking form went]
The symmetric linking form of Chapter~\ref{chap:ch06}, Rem.~5 does not enter Theorem~\ref{ch07:thm:lattice},
and the guess that a double-stranded count is a Gauss sum or an isotropic count on $\K(D_k)$
can be discarded: on $\K(D_{12})=\Z/3\oplus\Z/11$ any non-degenerate quadratic form has a
single isotropic vector and a Gauss sum of modulus $\sqrt{33}$, while $\#\DS(S,12)=10$. What
survives of the form is only the anti-invariance $f+\rho f=0$, which is the linear shadow of
$\rho$ being a duality rather than a symmetry.
\end{remark}

\begin{theorem}[Rank of the anti-invariant lattice]\label{ch07:thm:rank}
Let $\#\mathrm{fix}_V$ be the number of $\rho$-fixed vertices and $\#\mathrm{fix}_A$ the
number of $\rho$-fixed contents of $D_k(S)$. Then
\[
  \rank\Lm=\frac{|A|-|V|+2-\#\mathrm{fix}_A-\#\mathrm{fix}_V}{2},\qquad
  \rank\Lp=\frac{|A|-|V|-2+\#\mathrm{fix}_A+\#\mathrm{fix}_V}{2}+b_0 ,
\]
where $b_0\in\{1,2\}$ is the number of components. The expression does not depend on the
compaction: $|A|-|V|$ and $\#\mathrm{fix}_A-\#\mathrm{fix}_V$ are the same for the
unitig-contracted graph and for the full de Bruijn graph of $S\bar S$.
\end{theorem}

\begin{proof}
Let $C_1=\Z^A$, $C_0=\Z^V$ with the boundary $\partial$, and let $P_A,P_V$ be the
permutation matrices of $\rho$ on contents and vertices. By Lemma~\ref{ch07:lem:rho}(a),
$\partial P_A=-P_V\partial$, so $\varphi=(-P_A,P_V)$ is a chain endomorphism of the cellular
chain complex of $D_k$, and the Hopf trace formula gives
\[
  \tr(\varphi_*|H_1)-\tr(\varphi_*|H_0)=\tr(-P_A|C_1)-\tr(P_V|C_0)
  =-\#\mathrm{fix}_A-\#\mathrm{fix}_V .
\]
Here $H_1=Z_1$ with $\varphi_*=-\rho$, and $H_0=\Z^{b_0}$ with $\varphi_*$ the permutation of
components induced by $\rho$, whose trace is the number of $\rho$-fixed components: $1$ when
$b_0=1$, and $0$ when $b_0=2$ since then $\rho$ exchanges the halves
(Lemma~\ref{ch07:lem:rho}(c)). Hence $\tr(\rho|Z_1)=\#\mathrm{fix}_A+\#\mathrm{fix}_V-
\tr(\varphi_*|H_0)$. Over $\Q$ the involution $\rho$ is diagonalisable with eigenvalues
$\pm1$, so $\rank\Lp-\rank\Lm=\tr(\rho|Z_1)$ while $\rank\Lp+\rank\Lm=\rank Z_1=|A|-|V|+b_0$.
Solving,
\[
  \rank\Lm=\tfrac12\big(|A|-|V|+b_0+\tr(\varphi_*|H_0)-\#\mathrm{fix}_A-\#\mathrm{fix}_V\big),
\]
and $b_0+\tr(\varphi_*|H_0)=2$ in both cases. For the last sentence: unitig contraction is a
$\rho$-equivariant homotopy equivalence, so $\rank Z_1$ and the Lefschetz number
$\tr(\varphi_*|H_1)-\tr(\varphi_*|H_0)$ are unchanged, and the latter is
$-\#\mathrm{fix}_A-\#\mathrm{fix}_V$ in either cellular structure.
\end{proof}

\begin{corollary}\label{ch07:cor:parity}
For even $k$, $\rank\Lm=(|A|-|V|+2-\#\{\text{palindromic repeated $k$-mers}\})/2$; for odd
$k$, $\rank\Lm=(|A|-|V|+2-\#\{\text{palindromic $(k{+}1)$-mers of $S\bar S$}\})/2$. If $S$
has no inverted repeat of length $\ge k$ then $\Lm\cong Z_1(G_k(S))$.
\end{corollary}

\begin{proof}
The first two statements are Theorem~\ref{ch07:thm:rank} with Lemma~\ref{ch07:lem:rho}(b). If $D_k$
is disconnected then $Z_1(D_k)=Z_1(G_k(S))\oplus Z_1(G_k(\bar S))$ with $\rho$ exchanging the
summands, and $f\mapsto(f,-\rho f)$ identifies $Z_1(G_k(S))$ with $\Lm$.
\end{proof}

\begin{table}[t]
\centering
\begin{tabular}{rrrrrrrrr}
\toprule
$k$ & $|V|$ & $\#\mathrm{fix}_V$ & $|A|$ & $\#\mathrm{fix}_A$ & $\sum m$ & $b_0$ &
$\rank Z_1$ & $\rank\Lm$\\
\midrule
13 &   0 & 0 &   0 & 0 &   0 & -- & -- & 0\\
12 &   7 & 1 &  14 & 0 &  14 & 1 &   8 &   4\\
11 &  30 & 0 &  53 & 1 &  60 & 1 &  24 &  12\\
10 & 125 & 3 & 222 & 0 & 252 & 1 &  98 &  48\\
 9 & 396 & 0 & 685 & 3 & 812 & 1 & 290 & 144\\
\bottomrule
\end{tabular}
\caption{$\phi$X174 (NC\_001422.1, $5386$\,bp). $\rank\Lm$ is computed as the corank of the
balance matrix on orbit representatives, modulo two large primes, and agrees with
Theorem~\ref{ch07:thm:rank} in every row. At $k=13$ no $k$-mer is repeated on either strand and
$D_k$ is the pair of disjoint cycles $S,\bar S$. Check (L).}
\label{ch07:tab:rank}
\end{table}

\begin{remark}[Not the number of inverted repeats]\label{ch07:rem:notIR}
It is tempting to read $\rank\Lm$ as the number of independent inversions available to the
genome, and both the specification of this note and its planning did. Table~\ref{ch07:tab:rank}
refutes it. At $k=12$, $\phi$X174 carries exactly one inverted-repeat pair of length $\ge12$
(\texttt{ATTAAGCTCATT} at $441$ against its reverse complement at $1448$) and one palindromic
repeated $12$-mer, and $\rank\Lm=4$. At the other extreme, Corollary~\ref{ch07:cor:parity} says
that a genome with repeats but no inverted repeat has $\Lm$ as large as the whole circulation
lattice of $G_k(S)$, while only two of its box points --- $\spec(S)$ and $\spec(\bar S)$ ---
have connected support. The number that does count inversion classes is the number of
connected box points, $8$ and $218$ in Table~\ref{ch07:tab:lattice}, and it has no closed form we
know of.
\end{remark}

\begin{table}[t]
\centering
\small
\begin{tabular}{rrrrrrrr}
\toprule
$k$ & box & conn. & lattice & brute & un-norm. & one strand & $|\K(D_k)|$\\
\midrule
12 &  10 &   8 &  10 &  10 &   10 &  4 & 33\\
11 & 528 & 218 & 864 & 864 & 2756 & 66 & 25\,318\,912\\
10 & \multicolumn{5}{c}{rank $48$: beyond the enumerator} & $8\,610\,708\,632$ & $3.42\cdot10^{35}$\\
\bottomrule
\end{tabular}
\caption{The lattice sum for $\phi$X174. ``Box'' is $|\Lm\cap\{|f|\le m,\ f\equiv m\}|$;
``conn.'' the number of points with $\Nseq>0$; ``lattice'' is $\#\DS$ from
Theorem~\ref{ch07:thm:lattice}; ``brute'' the number of closed walks of $D_k$ with
$\mathrm{usage}+\rho\,\mathrm{usage}=m$, enumerated by backtracking without reference to the
lattice; ``un-norm.'' is $\sum_f t(c_f)\Loc(c_f)$; ``one strand'' is
$\Nseq(S)+\Nseq(\bar S)$. $\Loc(D_k)=1$ at $k=12,11$, so $|\K(D_k)|$ is also the Eulerian
circuit count of the double cover. Checks (B), (S), (E), (N).}
\label{ch07:tab:lattice}
\end{table}

\paragraph{What the lattice points are.}
At $k=12$ the eight connected points are $\spec(S)$ and $\spec(\bar S)$, each with
$\Nseq=2$ (the interleaved two-copy repeat of the single-strand theory), and six spectra with
$\Nseq=1$. One of the six is the inversion of $S[441..1459]$ between the inverted-repeat pair.
The other five mix that inversion with the repeat exchange and with hairpin turns at the
palindrome: for instance the point that drops the two segments $S[264..452]$ and
$S[1448..2771]$ of $S$ and takes their reverse complements is the sequence
\begin{align*}
  &S[441..457]\;\,S[3205..4483]\;\,S[4483..264{+}N]\;\,S[2760..3205]\;\,S[457..1448]\\
  &\qquad\overline{S[264..441]}\;\;\overline{S[1448..2760]},
\end{align*}
a rearrangement of the whole genome in which two blocks are read from the other strand. At
$k=11$ the $218$ connected points carry $\Nseq$ values $33$ (twice), $18,16$ (twice each),
$15$ (four times), $11$ (six), $10$ (two), $9$ (six), $6$ (22), $5$ (20), $4$ (six), $3$ (36),
$2$ (56) and $1$ (54), summing to $864$. Every one of the $874$ reconstructions at the two
levels is aperiodic, so $d(T)=1$ throughout and the weighted count is the count (check (A)).

\paragraph{The brute force.}
Check (E) enumerates, by depth-first search on the content multigraph, every closed walk
whose usage $u$ satisfies $u(a)+u(\bar a)\le m(a)$ along the way and $=m(a)$ at the end ---
the walk has exactly $\sum m/2$ arcs --- rooted at its smallest content $a_0$ and weighted
$1/u(a_0)$; the search never sees $\Lm$. It returns $90$ nodes and $8$ spectra at $k=12$,
$45\,061$ nodes and $218$ spectra at $k=11$; the spectra found are exactly the connected box
points, and the count at each spectrum is exactly $\Nseq$.

\section{The inversion flip}

\begin{definition}\label{ch07:def:flip}
Let $T$ be a double-stranded reconstruction, read as the closed walk $a_0a_1\cdots a_{L-1}$
of arcs of $D_k$ with vertex visits $v_i=\mathrm{tail}(a_i)$. A \emph{flip} at a pair of
visits $i\neq j$ with $v_j=\rho(v_i)$ replaces the block $a_i\cdots a_{j-1}$ (indices cyclic)
by $\bar a_{j-1}\cdots\bar a_i$:
\[
  T'=a_0\cdots a_{i-1}\;\bar a_{j-1}\cdots\bar a_{i}\;a_j\cdots a_{L-1}.
\]
The \emph{flip graph} has the reconstructions as vertices and the flips as edges.
\end{definition}

\begin{proposition}\label{ch07:prop:flip}
\begin{enumerate}
\item[\rm(a)] $T'$ is a double-stranded reconstruction, and every replacement of a block
$w\,X\,\bar w$ of $T$ by $w\,\overline X\,\bar w$, for a $k$-mer $w$ and a later occurrence of
$\bar w$, is a flip in the sense of Definition~\ref{ch07:def:flip}.
\item[\rm(b)] $\spec(T')=\spec(T)-c_B+\rho c_B$, where $c_B$ is the usage of the block. A flip
preserves the spectrum if and only if $c_B=\rho c_B$.
\item[\rm(c)] Flipping the complementary block gives $\overline{T'}$; flipping the whole
cycle at a palindromic visit gives $\bar T$.
\item[\rm(d)] If $S$ carries no inverted repeat of length $\ge k$, no reconstruction admits a
flip: the flip graph is discrete, with $2\Nseq(S)$ vertices.
\end{enumerate}
\end{proposition}

\begin{proof}
(a) The block $a_i\cdots a_{j-1}$ runs from $v_i$ to $v_j=\rho(v_i)$; its reverse-complement
image $\bar a_{j-1}\cdots\bar a_i$ runs from $\rho(v_j)=v_i$ to $\rho(v_i)=v_j$, so $T'$ is a
closed walk, and its spectrum is $\spec(T)-c_B+\rho c_B$, which still satisfies $c+\rho c=m$.
For the second statement: the block $w\,X\,\bar w$ has reverse complement
$\overline{\bar w}\,\overline X\,\bar w=w\,\overline X\,\bar w$, so the endpoints are
unchanged and the interior $(k{+}1)$-mers are replaced by their reverse complements, which
lie in $\spec(\bar T)$ and hence in the double-stranded spectrum. It remains to see that $w$
is a vertex of $D_k$: $w$ occurs in $T$ and $\bar w$ occurs in $T$, so $w$ occurs in $T$ and
in $\bar T$, so $w$ occurs at least twice in $T\bar T$, hence at least twice in $S\bar S$ by
the spectrum identity, hence is a repeated $k$-mer, i.e.\ a vertex. (b) is the computation in
(a). (c) The complementary block runs from $v_j$ to $v_i=\rho(v_j)$, and flipping it yields
the walk $\bar a_{i-1}\cdots\bar a_{0}\bar a_{L-1}\cdots\bar a_{j}\;a_j\cdots$, which as a
circular sequence is the reverse complement of $T'$; taking $i=j$ at a palindromic vertex
flips everything. (d) By Lemma~\ref{ch07:lem:rho}(c) $D_k$ is then disconnected, $T$ lies in one
half, and $\rho(v_i)$ lies in the other for every visit, so no pair of visits qualifies; the
reconstructions are those of $S$ and of $\bar S$.
\end{proof}

Part (d) says that connectivity of the flip graph is not a theorem: a genome with a
multi-copy repeat and no inverted repeat has $2\Nseq(S)>2$ reconstructions and no move between
any of them. What can be asked is whether a \emph{given} genome is flip-connected, and the
answer for $\phi$X174 is yes, in a strong form.

\begin{theorem}[$\phi$X174, computed]\label{ch07:thm:flips}
At $k=12$ and $k=11$:
\begin{enumerate}
\item[\rm(a)] the flip graph on all reconstructions ($10$, resp.\ $864$ vertices; $16$,
resp.\ $5140$ edges) is connected, without the move $T\mapsto\bar T$;
\item[\rm(b)] no flip preserves the spectrum: $c_B\neq\rho c_B$ for every block of every
reconstruction; consequently the quotient graph on spectrum classes ($8$, resp.\ $218$
vertices) is connected, and every class is at flip distance at most $3$, resp.\ $8$, from
$\spec(S)$ (histogram $1,2,3,2$, resp.\ $1,7,16,38,56,52,34,12,2$);
\item[\rm(c)] two reconstructions with the same spectrum are joined by flip paths of even
length at least four: at $k=11$ the $6776$ ordered same-spectrum pairs have flip distances
$4$ ($3932$ pairs), $6$ ($2556$) and $8$ ($288$); at $k=12$ all four have distance $4$;
\item[\rm(d)] the vertex pairs $\{v,\rho v\}$ realising flips number $3$, resp.\ $14$, of
which $1$, resp.\ $5$, are the inverted-repeat pairs and palindromes of $S$ itself; restricted
to those, flips reach $6$ of $8$, resp.\ $212$ of $218$, classes. The others are pairs
$\{a,\bar a\}$ with $a$ a repeated $k$-mer of one strand --- \texttt{AAACGGCAGAA} against
\texttt{TTCTGCCGTTT}, for instance --- which become flip sites only after an earlier flip has
brought $\bar a$ into the reconstruction.
\end{enumerate}
\end{theorem}

The theorem is the transcript of \texttt{ds\_flips.py}: all reconstructions are taken from the
brute force of check (E), all flips are generated from Definition~\ref{ch07:def:flip}, and each
flip is verified to land in the reconstruction set (13\,928 instances at $k=11$).

\begin{remark}[Transpositions as products of flips]
Part (c) is the de Bruijn-graph form of a classical fact. The two reconstructions of
$\spec(S)$ at $k=12$ differ by exchanging the two copies of a repeat, which as a rearrangement
of $S$ is a transposition of the blocks between the copies. In sorting by reversals
\cite{BP96,HP99} a transposition is a product of three reversals. Here a reversal is
available only when it is flanked by a vertex and its reverse complement, and the shortest
product that realises the exchange has four factors; the distances in (c) are all even, which
we have not explained. Part (d) shows the mechanism: the exchange needs a flip at
$\{a,\bar a\}$ with $a$ a single-strand repeat, and the only way to bring $\bar a$ into $T$ is
to flip first at an inverted-repeat pair of $S$, then at $\{a,\bar a\}$, and then to undo the
first flip.
\end{remark}

\begin{remark}[Corrections to the specification]\label{ch07:rem:corrections}
Three statements in the specification of this note, one of them originating in its planning,
are not correct as stated. \emph{First}, $\rank\Lm$ is not the number of independent inverted
repeats of length $\ge k$ (Remark~\ref{ch07:rem:notIR}); it is given by Theorem~\ref{ch07:thm:rank}.
\emph{Second}, the per-point term is not $t(c_f)\prod(d^+-1)!$ but that quantity divided by
$\prod_a c_f(a)!$ (Remark~\ref{ch07:rem:norm}); the un-normalised sum is $2756$ where the count is
$864$. \emph{Third}, ``the flip graph is connected'' and ``a same-spectrum transposition is a
composition of inversions'' are properties of $\phi$X174 at $k=12,11$, established by
computation, and not theorems: Proposition~\ref{ch07:prop:flip}(d) gives the general obstruction,
and the composition, where it exists, costs four flips rather than the three reversals of the
unconstrained model.
\end{remark}

\section{Scope, provenance, and what is open}

\paragraph{Scope.} Lemma~\ref{ch07:lem:rho}, Lemma~\ref{ch07:lem:spectrum}, Theorem~\ref{ch07:thm:lattice},
Theorem~\ref{ch07:thm:rank}, Corollary~\ref{ch07:cor:parity} and Proposition~\ref{ch07:prop:flip} are proved
in full generality for a single circular sequence. Theorem~\ref{ch07:thm:flips} is a computation
for one genome at two word lengths. The lattice sum is evaluated at $k=12$ and $k=11$; at
$k=10$ the lattice has rank $48$ and the capacity-pruned backtracking enumerator of
\texttt{ds\_lattice.py} produces no point in ten minutes, so only the rank row of
Table~\ref{ch07:tab:rank} is reported there. The count is a count of circular sequences, each
weighted by $1/d(T)$; in the data $d(T)=1$ throughout.

\paragraph{Provenance.} The BEST theorem is \cite{BEST,SmithTutte}; its use to count genome
reconstructions is \cite{KSP}; the bidirected-graph model of double-stranded assembly, in
which a double-stranded reconstruction is a walk in a bidirected de Bruijn graph, is
\cite{MGMB07} after \cite{KM}, and those papers ask for existence and for one solution, not
for the count. Sorting by reversals and the transposition-as-three-reversals identity are
\cite{BP96,HP99}. The double cover $D_k$, its balance and its connectivity criterion are
Chapter~\ref{chap:ch05}; the anti-automorphism theorem and the symmetric linking form are Chapter~\ref{chap:ch06}; the
single-strand counting convention is Chapter~\ref{chap:ch01}. The Hopf trace formula is standard. We are
not aware of prior work expressing a double-stranded reconstruction count as a sum over an
eigenlattice of the strand involution, or of prior work on the flip graph of a de Bruijn
graph; a targeted search found none.

\paragraph{Open.} \emph{First}, the number of connected box points --- the number of spectrum
classes, or equivalently the size of the orbit of $\spec(S)$ under the flip pseudo-group ---
has no formula; $8$ and $218$ are all that is known. \emph{Second}, the genome-scale question:
Theorem~\ref{ch07:thm:rank} evaluates the rank of $\Lm$ for \emph{E.\ coli} in one line from the
double cover of Chapter~\ref{chap:ch05}, but the sum is not enumerable there, and what is wanted is a
transfer-matrix or generating-function evaluation of $\sum_f\Nseq(c_f)$ along the genome, or
an Ehrhart-type asymptotic in $k$. \emph{Third}, the parity of same-spectrum flip distances in
Theorem~\ref{ch07:thm:flips}(c). \emph{Fourth}, the same lattice with the involution removed and
$m$ split into $p$ parts is the phasing problem for a polyploid, where the symmetric group on
the parts acts on the decompositions rather than on the graph; that is the subject of the next
note.
